\section{Apéndice: recopilación de datos y procesamiento de imágenes}
\subsection{Descarga de video y extracción de fotogramas}
Para obtener una portada nítida y los key frames, descargamos el video original con precisión de 720p mediante \texttt{yt-dlp -f 398+251 -o lecture11.webm} (enlace del video \texttt{https://www.youtube.com/watch?v=\_MNlLhU33H0}); luego, en los cuatro instantes 00:02:35, 00:36:00, 00:44:40 y 01:01:30, invocamos \texttt{ffmpeg -ss ... -frames:v 1 -update 1} para generar sucesivamente frame-01–frame-04. La tabla siguiente enumera los comandos clave y su propósito.

\begin{table}[H]
\centering
\begin{tabular}{p{6cm}p{9cm}}
\toprule
Comando & Descripción \\
\midrule
\texttt{yt-dlp -f 398+251 -o lecture11.webm https://www.youtube.com/watch?v=\_MNlLhU33H0} & Descarga el video en 720p+audio, y proporciona el material para las capturas y la portada posteriores. \\
\texttt{ffmpeg -ss 00:02:35 -i lecture11.webm -frames:v 1 -update 1 frame-01-selfreward.jpg} & Captura el key frame de la sección de self-rewarding. \\
\texttt{ffmpeg -ss 00:36:00 -i lecture11.webm -frames:v 1 -update 1 frame-02-iterative.jpg} & Captura la ilustración del iterative pipeline. \\
\texttt{ffmpeg -ss 00:44:40 -i lecture11.webm -frames:v 1 -update 1 frame-03-selfinstr.jpg} & Fotograma de ejemplo de Self-Instruct. \\
\texttt{ffmpeg -ss 01:01:30 -i lecture11.webm -frames:v 1 -update 1 frame-04-meta.jpg} & Imagen de Meta judge. \\
\bottomrule
\end{tabular}
\caption{Comandos clave de recopilación e imágenes.}
\end{table}

\subsection{Depuración de subtítulos y esquema}
Depuramos el SRT mediante Python: eliminamos los párrafos duplicados, conservamos las marcas de tiempo y extraemos la oración temática cada 2,500 líneas por bloque para apoyar la escritura. El núcleo del script divide los bloques con \texttt{re.split(r'\\n\\n+')}, extrae la marca de tiempo de la primera línea y usa \texttt{prev} para eliminar duplicados y conservar lo esencial. El archivo \texttt{subs\_clean.txt} resultante sirve, en la redacción de las notas, como fuente de la chronology; contiene palabras clave de párrafo como \texttt{Self-Rewarding} y \texttt{Meta Judge}, lo que facilita clasificarlo según la lógica didáctica.

\subsection{Resumen de la sección}
El proceso de descarga, captura de pantalla y procesamiento de subtítulos que documenta el apéndice constituye el respaldo de la escritura posterior, y garantiza que cada imagen y cada fragmento de texto puedan rastrearse hasta un instante y un script específicos.
